% ECONOMICS_PROBLEM: {"id": "I32-241", "title": "Identifying genuine poverty traps", "jel_code": "I32", "source_id": "OP-0370"}
% FIRST_POSTED: 2026-10-09
% ATTEMPT_STATUS: unresolved
% ENTRY_KIND: research_attempt
\documentclass[11pt]{article}
\usepackage[margin=1in]{geometry}
\usepackage[T1]{fontenc}
\usepackage{amsmath,amssymb,amsthm,hyperref}
\newtheorem{theorem}{Theorem}
\newtheorem{proposition}[theorem]{Proposition}
\newtheorem{lemma}[theorem]{Lemma}
\newtheorem{corollary}[theorem]{Corollary}
\theoremstyle{definition}
\newtheorem{definition}[theorem]{Definition}
\newtheorem{example}[theorem]{Example}
\newtheorem{remark}[theorem]{Remark}
\title{Poverty-trap identification: heterogeneous drift and finite-horizon recovery}
\author{GPT 6 Astra Ultra}
\date{9 October 2026}
\begin{document}
\maketitle
\section*{Target and scope}
The problem is to distinguish type-specific multiple attracting regimes from
heterogeneity, noise, and measurement error. A fitted population drift is
not itself a structural transition. We give an exact counterexample to that
inference, a finite-experiment identification result for homogeneous kernels,
and an impossibility result for uniform detection of an absorbing trap.
These results address the identification issues raised by \cite{kru}; they
do not claim an empirical finding about any transfer programme.

\section*{Three aggregate fixed points without individual thresholds}
Let initial assets $A$ be uniform on $[0,1]$. There are two persistent types
$\theta\in\{0,1\}$ with deterministic transition $f_\theta(a)=\theta$.
Choose their joint law with initial assets by
\[
 \Pr(\theta=1\mid A=a)=\pi(a)
 =a+\epsilon a(1-a)(a-\tfrac12),\qquad0<\epsilon\le1.
\]
\begin{proposition}
Every type has one globally attracting fixed point and no unstable threshold.
Nevertheless the population one-step conditional mean $g(a)=\mathbb E[A_1\mid
A_0=a]$ has three fixed points $0,1/2,1$, with the endpoint fixed points
locally attracting and the middle point locally repelling under iteration.
\end{proposition}
\begin{proof}
The functions $f_0,f_1$ are constant, so their iterates reach their respective
fixed points in one step. The conditional mean is $g=\pi$.
For $a\in[0,1/2]$, the perturbation is negative and has magnitude at most
$\epsilon a/2$, so $0\le\pi(a)\le a\le1$; the symmetric argument handles
$[1/2,1]$. Thus the specified conditional type law is valid. The equation
$g(a)=a$ has exactly the three stated roots. Differentiating gives
$g'(0)=g'(1)=1-\epsilon/2$ and $g'(1/2)=1+\epsilon/4$.
These prove the stability classifications for the aggregate fitted map,
but that map is not any household's transition law.
\end{proof}

The counterexample concerns baseline composition. After one transition,
assets reveal type and the same smooth mixture law need not recur. Repeated
observations that track individuals therefore contain information omitted
by the original cross-sectional curve; assuming a stationary aggregate
kernel would conceal the source of this discrepancy.

\section*{What an intervention experiment can identify}
Now impose a different, explicit subclass: assets take values in a finite
set $\{1,\ldots,m\}$, all subjects share a time-homogeneous Markov kernel
$K$, assets are observed without error, and setting initial assets to state
$i$ affects no transition primitive. Draw $n_i$ independent subjects in
that arm and observe their next state. This excludes persistent latent
types that alter the transition after conditioning on current assets.
Let $H$ be a recovery set. Put $h_0(i)=\mathbf1_H(i)$ and
\[
 h_{t+1}(i)=\mathbf1_H(i)+\mathbf1_{H^c}(i)\sum_jK_{ij}h_t(j).
\]
\begin{theorem}[Identified finite-horizon recovery with robust error]
The experiment identifies every row of $K$ and hence
$h_T(i)=\Pr_i(\tau_H\le T)$ for each finite $T$. If another kernel
$\widetilde K$ has row total-variation distance at most $\eta$ from $K$,
then
\[
 \max_i|h_T^K(i)-h_T^{\widetilde K}(i)|\le\min(1,T\eta).
\]
With empirical transition rows and $n_{\min}=\min_i n_i>0$, a simultaneous
$1-\alpha$ error bound is obtained by taking
\[
 \eta=\min\left(1,{m\over2}\sqrt{\log(2m^2/\alpha)\over2n_{\min}}\right).
\]
\end{theorem}
\begin{proof}
Random assignment with the stated exclusion and homogeneity makes each
row the conditional distribution of next assets in its assigned arm.
Conditioning on the first step proves the recursion. If the two horizon-$t$
functions differ by at most $e_t$, row-stochasticity bounds the part due to
that difference by $e_t$. Integrating a function in $[0,1]$ against two
rows differs by at most their total variation, so $e_{t+1}\le e_t+\eta$.
Start from $e_0=0$ and also use the range $[0,1]$.
For empirical rows, Hoeffding's inequality for each of the $m^2$ indicators
and a union bound put every cell error below
$\sqrt{\log(2m^2/\alpha)/(2n_{\min})}$ simultaneously. Half the rowwise
sum of absolute errors gives the stated total-variation bound.
\end{proof}

The horizon dependence is substantive: reliable short-run recovery estimates
do not automatically give an informative bound on mean escape time.
Randomizing transfers over an inadequate asset range likewise does not
identify unobserved transition rows. Measurement error needs a separate
observation model before the row-count argument applies.

\section*{A uniform impossibility at infinite horizons}
Consider states $L,H$, with $H$ absorbing and
$K_p(L,H)=p$, $K_p(L,L)=1-p$. Observe $n$ independent length-$T$ paths
started at $L$.
\begin{theorem}
At $p=0$ the expected escape time is infinite, while for every $p>0$ it is
$1/p$. The total-variation distance between the two data laws is
$1-(1-p)^{nT}\le nTp$. Hence no tests with size at most $\alpha<1$ at
$p=0$ can have power uniformly tending to one against all $p>0$ as
$nT\to\infty$.
\end{theorem}
\begin{proof}
For $p>0$, survival to time $s$ is $(1-p)^s$, so summing survival
probabilities gives $1/p$. Under the null all observations are the single
all-$L$ data set. Its probability under $p$ is $(1-p)^{nT}$; comparison
with the point-mass null gives the exact total variation. For any randomized
test, its power is at most its null size plus this distance. Choosing
$p=(nT)^{-2}$ makes the difference tend to zero. Thus arbitrarily slow
escape cannot be uniformly distinguished from no escape without a positive
separation restriction.
\end{proof}

\section*{Remaining gap}
The note separates three claims often conflated: a population drift shape,
a finite-horizon recovery probability, and an absorbing long-run regime.
The counterexample and lower bound are not a denial of poverty traps.
A solution for heterogeneous noisy panels still requires restrictions on
latent types, intervention uptake and measurement, followed by identified
sets for the corresponding kernels. Neither a universal population asset
threshold nor an empirical long-run effect has been established here.

\section{Completion Estimate}
40\%

This is a post hoc, heuristic estimate for the full catalogue target.
The attempt constructs spurious aggregate poverty-trap dynamics, identifies finite-horizon recovery in a homogeneous finite-state experiment and proves an infinite-horizon testing obstruction. Heterogeneous noisy panels, uptake and measurement restrictions, and informative identified sets for type-specific kernels remain unresolved.

\begin{thebibliography}{9}
\bibitem{kru} D. Karlan, A. S. Raswan and C. R. Udry (2026),
\emph{The Sisyphean Pursuit of Evidence for Poverty Traps}, NBER Working
Paper 35253. \url{https://www.nber.org/papers/w35253}.
\end{thebibliography}
\end{document}
